\section{결과와 분석}\label{sec:results}
\subsection{여러 기반 추정기에서의 보정 효과}
표~\ref{tab:backbones}는 같은 319장에 대한 각 기반의 Base와 N3를 비교한다. 코너 오차 중앙값은 YOLO에서 6.721에서 5.778픽셀, DOPE에서 12.570에서 7.469픽셀, ResNet에서 8.223에서 7.085픽셀로 감소하였다. \pckten 도 각각 63.425\%에서 68.587\%, 24.170\%에서 44.338\%, 52.261\%에서 57.223\%로 증가하였다. 동일 보정 절차를 기반별로 학습했을 때 초기 코너의 중심부 정확도와 전체 주석 기준 성공 코너 비율이 개선되었다.

위치·회전 중앙값의 점추정도 세 기반에서 감소하였다. 표~\ref{tab:pose_uncertainty_main}에서 YOLO와 DOPE의 두 중앙값 변화는 사후 세션 구간에서도 음수였다. 반면 ResNet의 위치 중앙값 차이 구간은 $[-2.077,0.202]$cm로 0을 포함한다. 세 기반의 점추정 감소를 모든 자세량의 확증적 개선으로 해석하지 않는다. 전체 P90·방향각 구간과 seed별 편차는 보충자료에 제시한다.
\input{tables/backbone_results}
\input{tables/pose_uncertainty_compact}

큰 오류에는 개선과 악화가 함께 존재한다. DOPE의 코너 P90은 51.282에서 53.570픽셀로 증가했고, 회전 P90 차이도 약 $+1.590$도였다. ResNet의 위치 P90은 약 $+2.759$cm 증가하였다. 각 기반·seed의 보정 전후 자세 산출 ID 집합은 같았으며 신규 자세 실패·복구는 0이었다. DOPE의 자세 산출은 210/319장으로, 실패 109장은 전체 분모에 남았다. 이는 N3가 결측이나 잘못된 검출을 새롭게 복구하는 방법이 아니라는 범위와 일치한다.

코너의 조건부 P90과 전체 실패 포함 통계도 다르다. YOLO의 조건부 P90은 43.890에서 42.134픽셀로 낮아졌지만, 실패 벌점 포함 전체 P90은 61.710에서 61.715픽셀로 거의 변하지 않았다. 따라서 관찰된 국소 보정 효과를 전체 극단 오류의 해결로 확대하지 않는다.

\subsection{치수와 회전 대응 감독의 추가 효과}
표~\ref{tab:ablation_results}는 동일한 YOLO의 통제 구성 N0--N3를 비교한다. 영상 기반 N0에 치수 분기를 추가한 N2의 코너 중앙값은 5.944에서 5.778픽셀로 감소했고, \pckten 은 약 1.054퍼센트포인트 증가하였다. 위치·회전 중앙값도 7.260cm·2.176도에서 7.011cm·2.090도로 낮아졌다. 다만 이것은 치수 분기 추가의 관측 효과이며, 정보 내용과 추가 파라미터의 독립 효과를 완전히 분리한 결과는 아니다.
\input{tables/ablation_results}

회전 대응만 추가한 N1은 N0보다 위치·회전 중앙값이 소폭 높았다. 치수 보정에 회전 대응을 추가한 N3도 N2와 코너 중앙값·\pckten 이 거의 같았으며, 위치 중앙값은 7.011에서 7.068cm로 증가하고 회전 중앙값은 2.090에서 2.070도로 감소하였다. 코너 P90은 약 0.326픽셀 낮아졌으나 변화 방향은 지표별로 달랐다. 대칭 감독의 일반적 필요성이나 전체 개선의 주원인으로 해석할 근거는 제한적이다.

실제 학습에서 원래 번호와 다른 전체 대칭 대응을 선택한 노출은 DOPE의 seed별 290·290·289회로 약 0.3\%였고, ResNet에서는 모두 0회였다. 합성 학습의 네 방향 회전 대응 행도 0개였다. 동일한 감독 절차를 적용했다는 사실과 그 절차가 목표를 실제로 바꾼 빈도를 구분해야 한다. 별도 정사각형 평가의 개선을 네 방향 대응 학습의 전이 효과로 설명하지 않는다.

\subsection{다른 보정 방법과의 비교}
표~\ref{tab:comparators}는 같은 초기 예측과 319장·8코너 평가에서 보정 방법들의 최종 결과를 비교한다. 제안 N3의 코너 중앙값은 직접 회귀 D와 선 구조 L보다 낮았다. 그러나 PoseFix 기반 방법은 코너 중앙값 5.561픽셀, 위치 중앙값 6.951cm, 회전 중앙값 2.028도로 N3보다 낮았고 \pckten 도 더 높았다. N3의 코너 P90은 42.134픽셀로 PoseFix 기반 방법의 43.907픽셀보다 낮았다.
\input{tables/comparator_results}

이 비교는 내부 특징·원영상 사용, 출력 방식, 손실과 감독 조건이 다른 전체 방법 비교이다. N3의 구조적 특징은 기존 특징과 치수를 이용해 제한된 좌표 이동을 학습한다는 데 있으나, 그것만으로 다른 보정 방법보다 정확하거나 효율적이라고 결론내리지 않는다. 동일 정보·예산의 완전 통제와 모든 비교군의 같은 환경 지연시간은 추가 확인이 필요하다.

\subsection{재질과 정사각형으로의 적용 범위}
직사각형 플라스틱 194장에서 코너 중앙값은 Base의 7.509에서 N3의 6.203픽셀로, 목재 125장에서는 6.124에서 5.274픽셀로 낮아졌다. 두 집단의 \pckten 과 위치·회전 중앙값도 감소 또는 증가의 개선 방향을 보였다. 목재는 재질 단위 평가가 완료된 것이며, 아직 확인되지 않은 가림 등급별 성능을 뜻하지 않는다. 집단의 시점·가림·촬영 조건이 달라 재질 자체의 인과 효과로 해석하지 않는다. 전체 재질별 행은 보충자료에 제시한다.

표~\ref{tab:square_results}의 정사각형 전체 수동 참조 모드에서 코너 중앙값은 YOLO 5.526에서 5.004픽셀, DOPE 11.748에서 7.037픽셀, ResNet 6.584에서 5.891픽셀로 낮아졌다. 영상 내부 600코너 모드에서도 각 기반의 보정 전후를 같은 분모로 제시한다. 두 모드 모두 2D 보정의 관측 근거이나, 한 촬영 세션·한 치수의 결과이므로 새로운 임의 개체·환경에 대한 일반화 범위는 제한된다.
\input{tables/square_results}

정사각형에서도 대칭 감독의 추가 이득은 확인되지 않았다. 전체 수동 모드에서 YOLO N3−N2는 코너 중앙값 약 $+0.055$픽셀, \pckten 약 $-0.332$퍼센트포인트였다. Base 대비 전체 보정의 이득과 N2 대비 대칭 감독의 추가 효과는 별개의 비교이다. 독립 표준 자세 참조가 없어 두 모드의 위치·회전 정확도는 모두 \notrun 이다.

\subsection{가림 조건과 큰 오류의 한계}
표~\ref{tab:occlusion_results}는 완료 제출된 사람 직접 분류를 원본 RGB 해시와 frame ID로 연결해 319장 전체를 clean 151장, moderate 87장, severe 81장으로 사후 재집계한다. clean과 moderate에서는 코너 중앙값·P90 및 위치·회전 중앙값이 Base보다 N3에서 낮았다. severe에서는 코너 중앙값이 11.084에서 9.696픽셀로 감소했지만 위치 중앙값은 13.484에서 15.706cm로 증가하였다. 회전 중앙값 61.635에서 11.706도로의 변화가 모든 프레임에서 같은 크기의 개선을 뜻하지는 않는다. 가설 선택 변화와 중앙 순위 이동을 포함한 프레임별 원인 확인이 필요하다.
\begin{table*}[t]
\centering\footnotesize
\setlength{\tabcolsep}{4pt}
\caption{사람 직접 분류를 완료한 319장의 가림 정도별 YOLO Base와 N3 사후 분석.}\label{tab:occlusion_results}
\begin{tabularx}{\textwidth}{Ylrrrrrrr}
\toprule
집단 & 방법 & 영상 & \shortstack{코너 중앙값\\(px)} & \shortstack{코너 P90\\(px)} & \shortstack{\pckten\\(\%)} & \shortstack{위치 중앙/P90\\(cm)} & \shortstack{회전 중앙/P90\\(도)} & \shortstack{참조/유효\\예측 코너} \\
\midrule
전체 & Base & 319 & 6.721 & 43.890 & 63.425 & 7.897 / 40.530 & 2.539 / 86.527 & 2,499 / 2,445 \\
전체 & N3: 제안 보정 & 319 & 5.778 & 42.134 & 68.587 & 7.068 / 37.809 & 2.070 / 85.919 & 2,499 / 2,445 \\
clean & Base & 151 & 5.272 & 19.351 & 73.960 & 4.464 / 18.976 & 1.716 / 10.802 & 1,202 / 1,202 \\
clean & N3: 제안 보정 & 151 & 4.525 & 15.575 & 80.616 & 4.131 / 15.388 & 1.477 / 5.847 & 1,202 / 1,202 \\
moderate & Base & 87 & 7.080 & 99.876 & 64.531 & 8.167 / 29.192 & 2.296 / 83.774 & 671 / 671 \\
moderate & N3: 제안 보정 & 87 & 5.918 & 99.665 & 67.263 & 7.215 / 27.063 & 1.944 / 83.139 & 671 / 671 \\
severe & Base & 81 & 11.084 & 52.325 & 42.013 & 13.484 / 173.736 & 61.635 / 89.349 & 626 / 572 \\
severe & N3: 제안 보정 & 81 & 9.696 & 51.516 & 46.912 & 15.706 / 175.733 & 11.706 / 89.041 & 626 / 572 \\
\bottomrule
\end{tabularx}
\tabnote{모든 행의 자세 산출률은 100\%이며 정답률이 아니다. N3는 세 seed의 통계 평균이다. 등급은 2026-09-22 제출된 사람 직접 분류이며, 모델 오차로 만든 레이블이 아니다. 이 표는 개발자료의 사후 분석이다. P·N2의 같은 집단 결과와 전체 분모는 보충자료에 제시한다.}
\end{table*}


새 등급은 플라스틱 194장과 목재 125장을 모두 포함하며 미분류는 0장이다. 새 레이블로 집단만 바꾸었을 때 전체 319장의 기존 2D·6D 수치가 정확히 유지됨을 확인했다. 다만 이 분류와 모델 선택에 같은 개발자료가 사용되었으므로 독립 시험으로 해석하지 않는다. 코너 가시성의 보조 분석에서는 직접 가시 66점의 중앙값은 낮아졌지만, 외부 가림 5점의 \pckten 은 두 방법 모두 0이었다. 나머지 참조 코너 2,428점은 가시성 사람 검수 전이므로, 작은 확정 표본으로 비가시 코너 복원의 일반적 성능을 판단할 수 없다.

초기 코너 오차를 $e_k$, 원본 영상에서의 이동 상한을 $c_i=0.01D_{I_i}$라 하면, 고정 참조 대응에 대해 다음 하한이 성립한다.
\begin{equation}
 e'_k\geq\max(0,e_k-c_i).
\end{equation}
640$\times$480 영상에서는 상한이 8픽셀이므로 초기 오차 50픽셀을 한 번에 42픽셀보다 작게 만들 수 없다. 이는 제한된 이동의 기하적 성질이지, 현재 모든 큰 오류의 원인이 이동 제한이라는 실험적 결론은 아니다.

표~\ref{tab:tail}에서 초기 오류가 상한 밖인 코너는 981/2,445개였다. N3가 20픽셀 초과 오류를 10픽셀 미만으로 복구한 코너 수는 seed별 0·2·0개였고, 5픽셀 미만 오류를 10픽셀 초과로 악화시킨 수는 0·0·1개였다. 또한 코너가 좋아져도 위치 또는 회전이 악화되는 영상이 남았다. 최종 효과는 큰 오류의 전반적인 복구보다 국소 정밀도 개선에 가깝다.
\input{tables/tail_results}

그림~\ref{fig:examples}는 코너가 참조에 가까워지는 사례와 처음에는 정확했던 점이 소폭 악화되는 사례를 함께 보여준다. 표시된 사례 오차는 해당 영상의 평균 코너 오차이며, 주표의 전체 코너 풀링 중앙값과 다르다. 큰 오차가 27.21에서 25.42픽셀로 줄어든 예도 완전한 복구가 아니라 부분적인 감소이다. 이 예시를 새 정사각형 또는 가림 등급 전체를 대표하는 확률 표본으로 해석하지 않는다.
\input{figures/examples_panel}

\subsection{추가 비용과 추정기 업데이트 대안}
표~\ref{tab:cost}에서 YOLO의 전체 CUDA 중앙 시간은 11.195에서 15.138ms로 증가하였다. 보정 단독 중앙 시간은 3.432ms이며 전체 중앙값의 차이 3.943ms와 다른 계측량이다. DOPE는 62.736에서 65.807ms, ResNet은 8.890에서 11.592ms로 증가하였다. 같은 RTX 3080에서 라이브러리 환경을 구분해 보고하며, 추가 비용을 실제로 측정했다는 범위로 해석한다. 이 결과가 비교 보정기 전반에 대한 효율 우월성이나 Jetson의 실행시간을 보장하지는 않는다.
\input{tables/cost_results}

표~\ref{tab:update_alternative}는 공통 플라스틱 128장으로 제한한 추정기 업데이트 대안이다. N3의 \pckten 은 56.244\%로 보정 의사 레이블 학생의 51.472\%보다 높았으며 위치·회전 중앙값도 낮았다. 합성 추가 학습과 원래 의사 레이블 학생을 포함한 대안 간 차이는 이 집합과 각 감독 예산에 한정된다. 별도 교사를 사용한 학생을 N3를 학습에 적용한 방법으로 명명하지 않으며, 자기학습 전체가 열등하거나 과적합이 실패 원인으로 확정됐다고 해석하지 않는다.
\input{tables/update_results}
